\subsection{Contexto}

Los modelos de lenguaje de gran escala (LLMs, por sus siglas en inglés) han alcanzado
un nivel de capacidad sin precedentes en los últimos años. Modelos como GPT-5
\citep{openai2025gpt5}, Claude Opus 4.6 \citep{anthropic2025opus} y Gemini 3
\citep{google2025gemini3} demuestran desempeño competitivo en razonamiento
matemático, generación de código, análisis de texto y resolución de problemas
complejos. En múltiples escalas académicos y profesionales, estos sistemas
igualan o superan el rendimiento humano en tareas que hace pocos años se
consideraban fuera del alcance de la inteligencia artificial.

Sin embargo, conforme las capacidades fundamentales de los LLMs se consolidan,
la frontera de investigación y desarrollo se ha desplazado hacia un problema
distinto: la \textbf{gestión eficiente del contexto}. La ventana de contexto
---el espacio limitado de tokens que un modelo puede procesar en una sola
interacción--- se ha convertido en el recurso más valioso y escaso en
aplicaciones para LLM's. Cada token consumido representa un costo económico
directo, una contribución a la latencia de respuesta y una porción del espacio
disponible para información relevante. En este sentido, el desafío ya no es si
un LLM \textit{puede} resolver una tarea, sino \textit{cuán eficientemente}
se puede minimizar el gasto al resolver dicha tarea.

Este problema se intensifica con la adopción de \textbf{sistemas agénticos}, en
los cuales un LLM no opera de forma aislada sino que interactúa con herramientas
externas ---bases de datos, APIs, sistemas de archivos, servicios web--- para
completar tareas del mundo real. El atributo dominante para esta interacción es
el \textit{tool calling}: el modelo recibe las definiciones (schemas) de las
herramientas disponibles dentro de su contexto, selecciona cuál invocar, genera
los parámetros apropiados para dicha tarea, recibe el resultado, y continúa su razonamiento.
Cada ciclo de este proceso consume tokens: las descripciones de herramientas,
los parámetros serializados, las respuestas intermedias y el historial acumulado
de la conversación compiten por espacio dentro de la ventana de contexto.

En 2024, Anthropic introdujo el \textbf{Model Context Protocol} (MCP), un
protocolo abierto que estandariza la comunicación entre LLMs y servidores de
herramientas \citep{anthropic2024mcp}. MCP define una interfaz cliente-servidor
basada en JSON-RPC que permite a los modelos descubrir e invocar herramientas de
manera uniforme, independientemente de su implementación subyacente
\citep{mcpspec2024}. La adopción del protocolo ha crecido rápidamente en la
industria, estableciéndose como un estándar de facto para la integración de
herramientas en sistemas basados en LLMs.

Dentro del ecosistema MCP, han emergido dos paradigmas fundamentalmente
distintos para la ejecución de tareas que requieren herramientas:

\begin{enumerate}
    \item \textbf{Llamadas directas a herramientas}: el modelo invoca
    herramientas individuales a través de MCP, una por una, acumulando
    resultados en su contexto a lo largo de múltiples turnos de interacción.

    \item \textbf{Ejecución de código}: el modelo genera un programa
    ---típicamente en Python--- que orquesta múltiples llamadas a herramientas
    dentro de un entorno de ejecución aislado (\textit{sandbox}), retornando
    únicamente el resultado final al contexto del modelo
    \citep{anthropic2024codeexec}.
\end{enumerate}

Ambos paradigmas resuelven el mismo problema funcional ---conectar un LLM con
herramientas externas--- pero difieren radicalmente en cómo utilizan la ventana
de contexto. Las llamadas directas cargan todas las definiciones de herramientas
y resultados intermedios en el contexto del modelo; la ejecución de código
delega la orquestación a un programa externo, reduciendo potencialmente el
número de tokens consumidos a cambio de un overhead inicial de generación de
código y ejecución en sandbox.

Esta distinción tiene implicaciones directas en las tres dimensiones que
determinan la viabilidad de un sistema agéntico en producción: el
\textbf{consumo de tokens} (costo económico), la \textbf{latencia} (experiencia
del usuario) y la \textbf{precisión} (confiabilidad del sistema). No obstante,
hasta la fecha no existe un estudio sistemático que cuantifique estos trade-offs
de manera controlada y reproducible en sistemas agenticos. El presente trabajo busca llenar este vacío.
